\section{The coefficient-prime tame character and the ordered ordinary flag}
\label{sec:cubic-coefficient-tame-refinement}

This section concerns the actual coefficient prime \(p=269\).
It refines AC-G78 and AC-G78-type without changing their
hypotheses on the inherited odd-\(e\) Type II endpoint, their
host, or their qualitative scope.  In particular
\(F_{\widetilde v}=K_{269}\) at each chosen coefficient-prime
place, an unramified quadratic extension of \(\mathbf Q_{269}\).
The character \(\chi\) has nontrivial tame cubic inertia here.
We keep the geometric-Frobenius convention of AC-G78.

\begin{lemma}[AC-G79-torus: an integral full-unit character summand]
\label{lem:G79-torus}
Let
\[
 T_p^0=\prod_{v\mid p}T(\mathcal O_{F_{\widetilde v}}),
 \qquad
 \Delta_p=\prod_{v\mid p}T(k_{\widetilde v}),
 \qquad
 T_p^0=\Delta_p\times\mathfrak t_p,
 \tag{G79.1}
\]
where the last equality uses the multiplicative Teichmuller
section.  After a finite unramified coefficient extension,
every character \(\delta:\Delta_p\to\mathcal O^\times\)
defines an integral idempotent
\[
 e_\delta=\frac1{|\Delta_p|}
                 \sum_{a\in\Delta_p}\delta(a)^{-1}R_a.
 \tag{G79.2}
\]
The entire construction of AC-G78-module restricts to this
summand, giving a finite free
\(\mathcal O[[\mathfrak t_p]]\)-module
\(\mathscr M_\delta=e_\delta\mathscr M\) and the actual
specialization
\[
 \mathscr M_\delta\otimes_{\mathcal O[[\mathfrak t_p]],
                         w_0\lambda^{-1}}\mathcal O'
       \simeq e_\delta H_\lambda(K^p;\mathcal O')
 \tag{G79.3}
\]
for every continuous weight and finite coefficient extension.
All maps here commute with the away Hecke operators.
The inverse-transpose map of AC-G75 sends this summand to
the summand with character
\[
                   \delta^\Theta(a)=\delta(\Theta(a))
                                    =w_0\delta^{-1}(a).
 \tag{G79.4}
\]
The four elliptic tests of AC-G75, and its conclusion from an
actual full-chain characteristic-zero fiber, remain valid
in these summands.
\end{lemma}

\begin{proof}
The order of \(\Delta_p\) is prime to \(p\), so (G79.2) is
integral.  An element of \(T_p^0\) normalizes \(N(\mathcal O_p)\)
and its conjugate by
\(\eta=\operatorname{diag}(\varpi^2,\varpi,1)\), and commutes
with \(\eta\).  Conjugation therefore permutes the cosets in
the literal scalar \(U_p\)-sum in Thorne's Section~6.
Thus the full unit torus, not only \(\mathfrak t_p\), commutes
with \(U_p\), the integral ordinary projector, all depth
inclusions, and the trace maps.  It also commutes with every
away Hecke correspondence.

At each finite depth the \(\delta\)-summand of the ordinary
module in AC-G78-module is a projective direct summand over
the local group ring
\(\mathcal O[\mathfrak t_p/\mathfrak t_p(d)]\), hence is free.
The trace isomorphisms and the identity
\[
 I_{d,d+1}(\operatorname{Tr}m)
       =\sum_{a\in\ker(D_{d+1}\to D_d)}a\,m
\]
restrict to this summand.  Lifting bases and applying the
twisted norm on each regular group-ring factor proves
(G79.3), with exactly the trace transitions and arbitrary
continuous-character specialization of AC-G78-module.
Equivalently \(e_\delta\mathscr M\) is finite projective over
the local Iwasawa algebra and therefore free.

For the involution, \(C R_a C=R_{\Theta(a)}\); changing the
index in (G79.2) gives (G79.4).  The identity
\(C U_p C=R_zU_p\) of AC-G75 still holds, and the central
translations \(R_z^{\pm1}\) commute with these projectors.
Thus \(C\) is an isomorphism between the specified integral
ordinary summands, with the transformed weight, tame level,
and Hecke localization stated in AC-G75.

Finally the finite-torus action is at \(p\), while the
elliptic test complex and all its group-algebra actions in
AC-G75 are at \(v_0\nmid p\).  It commutes with every
differential and comparison map; at a classical weight the
finite-torus action is transported by lowest-weight evaluation
as in (G79.9) below.  Applying the integral
idempotent therefore gives a direct summand of each complex.
Concentration and derived tensor identities restrict to
this summand, as do the finite free control and Nakayama
arguments.  If there is no Steinberg constituent in the
original \(\delta\)-summand at the fixed classical weight,
there is none in the transformed
\(\delta^\Theta\)-summand.  The three original tests and the
transformed fourth test consequently give the same
contradiction to the actual full-chain fiber as in AC-G75.
This proves the claimed restriction; it introduces no
inference from residual Satake congruence to exact fiber.
\end{proof}

\begin{lemma}[AC-G79-flag: the ordinary character and its Galois order]
\label{lem:G79-flag}
Let \(\sigma\) be an automorphic constituent on the definite
group in AC-G78-type, of parallel weight
\(\mu(e)=(e-1,0,1-e)\), ordinary in the source sense.
At a chosen coefficient-prime place, let
\(\alpha_1,\alpha_2,\alpha_3\) be the inducing characters in
increasing \(p\)-adic valuation order.  The scalar
continuous-function ordinary line has full unit-torus
character
\[
       \delta_{\rm full}
          =(\alpha_1\otimes\alpha_2\otimes\alpha_3)
                              \,w_0(\mu(e))^{-1}.
 \tag{G79.5}
\]
Suppose its pro-\(p\) character is \(w_0(\mu(e))^{-1}\)
and its Teichmuller character is the restriction of
\[
       (\xi_v,1,\xi_v)\,w_0(\mu(e))^{-1},
 \qquad
       \xi_v\ \longleftrightarrow\
                 \chi^j\tau|_{G_{F_{\widetilde v}}},
 \tag{G79.6}
\]
where the arrow is local reciprocity in the convention of
Thorne, and \(\tau\) is the trivial or quadratic conductor-\(269\)
character of AC-G78.
Then the smooth inducing characters have the exact ordered
inertial restrictions
\[
                (\alpha_1,\alpha_2,\alpha_3)|_{\mathcal O^\times}
                           =(\xi_v,1,\xi_v).           \tag{G79.7}
\]
Let \(r=\epsilon r_{\sigma,\iota}\), the unitary rank-three
normalization of AC-G78-type.  Its restriction to
\(G_{F_{\widetilde v}}\) is potentially crystalline, has
zero monodromy, and has a full invariant flag whose ordered
graded characters are
\[
 u_1\,\chi^j\tau\,\epsilon^e,\qquad
 u_2,\qquad
 u_3\,\chi^j\tau\,\epsilon^{-e},
 \tag{G79.8}
\]
where each \(u_i\) is unramified.  These are actual
characteristic-zero characters; no assertion that an
unramified value \(u_i(\operatorname{Frob})\) equals \(1\)
is made.
\end{lemma}

\begin{proof}
We spell out the full-unit normalization, rather than
deducing it from a pro-\(p\) character.
In the source algebraic coefficient model, the combined
torus action is
\[
                (t\cdot f)(g)=\rho_\mu(t)f(gt).
\]
Choose the lowest-weight functional \(\ell\), so that
\(\ell(\rho_\mu(t)v)=w_0(\mu)(t)\ell(v)\) for every
\(t\in T(\mathcal O)\).
For \(F=\ell\circ f\), this gives
\[
        R_tF=w_0(\mu)(t)^{-1}\ell(t\cdot f).
 \tag{G79.9}
\]
The canonical map from the simultaneous invertible-\(U\)
part to the normalized Jacquet module is equivariant for
the entire torus.  The relevant primary proof is
Clozel--Thorne, \emph{Level raising and symmetric power
functoriality, I}, Lemma~2.5, printed pp.~4--5:
the canonical projection identifies that part with the
Jacquet invariants, and the unique ordinary tuple is
\(\alpha_1\otimes\alpha_2\otimes\alpha_3\) in increasing
valuation order.  Its proof also gives
\[
 v_p(\alpha_1(\varpi)),\
 v_p(\alpha_2(\varpi)),\
 v_p(\alpha_3(\varpi))
                  =(-fe,0,fe),                        \tag{G79.10}
\]
where \(f=[F_{\widetilde v}:\mathbf Q_p]=2\).
Thorne's Proposition~6.5 uses exactly this lemma.

One can check the common normalization directly.
Here \([N(\mathcal O):\eta N(\mathcal O)\eta^{-1}]=Q^4\)
and \(\delta_B(\eta)^{1/2}=Q^{-2}\), with \(Q=p^f\).
The literal scalar \(U_p\) therefore has, on the ordered
Jacquet line after lowest-weight evaluation, eigenvalue
\[
          Q^2\,\alpha(\eta)\,w_0(\mu)(\eta)^{-1}.
\]
Its valuation is
\(2f-2fe+2f(e-1)=0\).
Every other permutation of the three distinct slopes in
(G79.10) has strictly greater valuation, by the strict
rearrangement inequality for the coefficients \(2,1,0\).
Equivalently this is the product of the two normalized
operators in the proof of Clozel--Thorne's lemma.
Their integrality on the classical lattice and a unit
product force the individual ordinary eigenvalues to be
units.  The torus modulus is \(1\) on units, so (G79.9)
gives (G79.5) with no additional unit character.
Combining the pro-\(p\) character with the specified
Teichmuller character proves (G79.7).

We next verify the Galois assertion, including \(e=1\).
Thorne's Corollary~4.3(2)--(3), printed p.~21, supplies de
Rhamness and the full Frobenius-semisimplified local--global
compatibility at every split place, including \(p\).
After the common twist \(r=\epsilon r_{\sigma,\iota}\),
the relevant smooth parameter is the unitary one used in
(G79.10).  Its ordered finite inertial characters are
\(\chi^j\tau,1,\chi^j\tau\).
The first character is nontrivial: its cubic part is
nontrivial and cannot be cancelled by the quadratic \(\tau\).

Monodromy commutes with inertia, and a nonzero entry of
monodromy connects Frobenius slopes differing by \(f\).
The two equal finite-character endpoint lines have slope
difference \(2ef\), which is never \(f\) for integral
\(e\ge1\).  The middle line has a different finite inertia
character from each endpoint.
Hence every entry of monodromy is zero.
In particular when \(e=1\), the adjacent slope differences
are \(f\), but their different inertia characters still
exclude those entries.  This explicitly uses the ramified
cubic character at \(p=269\).
De Rhamness and the \(p\)-adic monodromy theorem now imply
potential crystallinity.

For clarity we give the filtered-module argument which
fixes the order of the Galois flag over the unramified
quadratic local field.  Choose a finite Galois extension
\(L/F_{\widetilde v}\) over which \(r\) is crystalline.
Enlarge the coefficient field if necessary.
The normalized Newton slopes are \(-e,0,e\), and at every
embedding the Hodge filtration jumps are \(-e,0,e\), in
the convention in which \(\epsilon\) has filtration jump
\(-1\).
More explicitly, Clozel--Thorne I, Theorem~2.2, gives
\(\{1-e,1,1+e\}\) for \(r_{\sigma,\iota}\); the twist
by \(\epsilon\) subtracts \(1\), giving the jumps just stated.

In \(D_{\rm cris,L}(r)\), let \(D_1\subset D_2\) be the
sums of the lowest one and two Newton-slope subspaces.
They are the unique such slope subobjects.
Use Newton and Hodge numbers normalized by the number
of embeddings of \(L\), so that parallel weights give
these same numbers.  For any \(k\)-dimensional filtered
subspace, its induced Hodge jumps have sum at least the
sum of the \(k\) smallest ambient jumps.
For \(D_k\), this lower bound equals \(t_N(D_k)\).
Weak admissibility gives the reverse inequality
\(t_H(D_k)\le t_N(D_k)\).  Thus equality holds, and
equality holds at each embedding, since each embedding
has the same individual lower bound.
Every further subobject of \(D_k\) inherits the weak
admissibility inequality from the ambient module.
Consequently \(D_1,D_2\) are weakly admissible.
The same conclusion holds after forgetting the coefficient field:
all further rational subobjects inherit the ambient inequality,
and the equality of total numbers scales by the coefficient degree.
The coefficient action is recovered by full faithfulness.

The Colmez--Fontaine equivalence produces invariant
subrepresentations over \(L\).
The unique Newton-slope characterization makes these
subobjects stable under the descent data, so the flag
descends to \(F_{\widetilde v}\).
The ordered filtration jumps of its rank-one quotients
are \(-e,0,e\); their finite inertia is, in the same order,
\(\chi^j\tau,1,\chi^j\tau\).
The rank-one crystalline equivalence, or local class
field theory after removing these parallel cyclotomic
powers and finite characters, identifies the quotients
as (G79.8): a rank-one de Rham character of weight zero
and trivial Weil--Deligne inertia is unramified.
This proves the assertion without applying a theorem
whose local field is restricted to \(\mathbf Q_p\).
\end{proof}

\begin{theorem}[AC-G79: the actual coefficient-prime refinement]
\label{thm:G79}
Retain every hypothesis of AC-G78-type, in particular
the actual coefficient prime \(269\) and odd \(e\).
Its automorphic representation \(\sigma\) can be chosen
so that, simultaneously at every chosen place over
\(269\), (G79.7) and (G79.8) hold.
It therefore has the same ordered finite inertia and
cyclotomic Hodge assignment as the inherited cubic
endpoint--middle--endpoint datum.
Its \(q\)-component still has the finite type
\(\tau_\theta\) of AC-G78-type, its global normalized
rank-three Galois representation is irreducible, and
the symmetric classical weight is unchanged.

The assertion does not impose the full original
away-place local types, prove descent to the original
tame level, bound \(q\), or establish per-curve containment
of the actual modified deformation ideal.
A subsequent application of the cover in AC-G73 still
requires its own splitting, host, and local-type hypotheses.
\end{theorem}

\begin{proof}
For each coefficient-prime place let \(\xi_v\) be as in
(G79.6), and define the fixed character of \(\Delta_p\)
by
\[
 \delta_e=
   \left.
    \left(\bigotimes_{v\mid p}(\xi_v,1,\xi_v)\right)
                    w_0(\mu(e))^{-1}
   \right|_{\Delta_p}.                                \tag{G79.11}
\]
All values lie in a finite unramified extension of the
coefficient field and have order prime to \(269\).

In the construction of AC-G78, take the same dense
classical progression \(w\equiv e\pmod{p-1}\).
At every such point, the endoscopic local inducing
characters have the same tame unit parts
\((\xi_v,1,\xi_v)\): the ordinary underlying form is
unramified at \(p\), while \(\tau\) and \(\chi^j\) are
fixed.
The parallel algebraic characters
\(\mu(w)\) and \(\mu(e)\) have the same restrictions
to \(\Delta_p\).  Indeed those restrictions factor
through the residue-field norm to \(\mathbf F_p^\times\),
and the exponent difference is divisible by \(p-1\).
Thus (G79.5) puts every dense classical endoscopic
point in the single summand \(e_{\delta_e}\).

Apply AC-G79-torus and repeat the actual finite-Hecke-image
argument of AC-G78 in \(\mathscr M_{\delta_e}\).
This is a free Iwasawa module, not a merely pointwise
subspace.  All dense classical trace identities therefore
give the same homomorphism from its actual Hecke image
to the Hida branch \(R\).
The analytic root and the rational prime \(q\) are
unchanged.  The arbitrary-base-change Cayley--Hamilton
argument again kills the kernel of the specialized
Hecke-image map, so the exact full-chain point occurs
in \(e_{\delta_e}H_{\lambda_q}\).
There is no additional density assertion at that
nonalgebraic weight.

The refined four-test argument of AC-G79-torus now gives
a Steinberg constituent at the original classical weight
inside \(e_{\delta_e}H_{\mu(e)}\).
The transformed character used for the fourth test is
\(\delta_e^\Theta\), as in (G79.4); this is precisely the
transformed summand under \(C\), and is not silently
identified with \(\delta_e\).

Next repeat AC-G78-type on the classical lattice with
this additional commuting finite-torus projector.
On the algebraic coefficient model this is the projector
transported by lowest-weight evaluation; it selects
the smooth character \((\xi_v,1,\xi_v)\), whereas on the
scalar function model its character is \(\delta_e\).
These two character conventions are related by (G79.9).
The projector is integral and commutes with
\(\mathcal O[\Gamma_q]\), the ordinary projector and
the residual away-Hecke idempotent.
Hence the resulting actual lattice is still projective
over \(\mathcal O[\Gamma_q]\).

The refined Steinberg constituent gives a positive
finite Steinberg multiplicity difference in this
lattice.  Every other constituent remains generic at
all split places by the same residual genericity witness.
Consequently the nonnegative-contribution argument and
the projective Brauer identity of AC-G78-type give an
actual \(\tau_\theta\)-constituent in this very summand.
Its \(q\)-local parameter and global irreducibility are
then exactly those already proved there; neither
argument discards the coefficient-prime projector.

For this constituent the pro-\(p\) torus character is
\(w_0(\mu(e))^{-1}\), by the definition of
\(H_{\mu(e)}\), and its finite character is (G79.11).
AC-G79-flag gives both (G79.7) and the actual ordered
Galois flag (G79.8), with zero monodromy.
The three requirements are therefore simultaneous:
the genuine finite type at \(q\), the prescribed
coefficient-prime ordered inertia and Hodge assignment,
and global rank-three irreducibility.
The remaining assertions in the statement delimit the
unchanged quantitative and deformation-support questions.
\end{proof}

\begin{remark}[Source and geometric scope]
The new finite-idempotent and filtered-subobject arguments
are proved above.  The exact primary interfaces used are
Thorne, Sections~6.1--6.2 and Proposition~6.5, including
its full classical coefficient convention; Corollary~4.3(2)--(3);
and Clozel--Thorne I, Definition~2.4 and the complete
printed proof of Lemma~2.5, pp.~4--5.
Colmez--Fontaine, \emph{Construction des representations
\(p\)-adiques semi-stables}, Theorem~A and Proposition~4.2, is a retained
foundational equivalence, not a new proof here.
No complete upstream ordinary-control, local--global
compatibility, or \(p\)-adic monodromy theorem is claimed
to have been reproved.

In geometric language the finite projector specifies a
character of a prime-to-\(269\) finite torus cover in the
ordinary tower.  The Galois flag is a flag of the actual
local system after the stated finite base change.
Neither statement identifies the support of an arbitrary
characteristic-\(269\) deformation curve with automorphic
Hecke support.
\end{remark}
